% Vietnamese translation for OI Public Library
% Source-PDF: IOI中国国家候选队论文/国家集训队2024论文集.pdf
\documentclass[11pt]{article}
\usepackage{oipl-vn}

\title{Tập luận văn đội tuyển quốc gia Trung Quốc năm 2024}
\author{Tập luận văn đội tuyển quốc gia Olympic Tin học}
\date{China Computer Federation, 2024}

\begin{document}
\maketitle

\origtitle{Tập luận văn đội tuyển quốc gia Trung Quốc Olympic Tin học năm 2024}
\sourcepath{IOI China candidate team papers / National training team 2024 collection.pdf}

\renewcommand{\abstractname}{Tóm tắt}
\begin{abstract}
PDF nguồn là tập hợp luận văn đội tuyển quốc gia Trung Quốc năm 2024. Tệp được
dàn trang bằng Adobe InDesign; phần bìa và mục lục có lớp văn bản đọc được,
nhưng nhiều trang thân bài trích xuất thành mojibake và sinh cảnh báo phông
chữ. Bản LaTeX này chuyển ngữ phần nhận diện tập sách và toàn bộ mục lục, đồng
thời dịch bài đầu tiên của Zhou Kangyang về bài toán tính tổng hàm nhân tính
dựa trên OCR và đối chiếu ảnh trang PDF.
\end{abstract}

\section{Thông tin đầu sách}

Tập sách có tiêu đề \emph{Tập luận văn đội tuyển quốc gia Trung Quốc Olympic
Tin học năm 2024}. Tệp PDF gồm 424 trang và được tạo bằng Adobe InDesign. Phần
văn bản trích xuất được không ghi riêng tên huấn luyện viên.

\section{Mục lục đã dịch}

\begin{center}
\small
\begin{tabular}{r p{0.52\linewidth} p{0.24\linewidth}}
\textbf{Trang} & \textbf{Bài viết} & \textbf{Tác giả} \\
\hline
1 & Một số tiến triển về bài toán tính tổng hàm nhân tính & Zhou Kangyang \\
11 & Bàn về bài toán parity của matroid tuyến tính & Guo Yuchong \\
24 & Bàn về hợp nhất và tách cấu trúc dữ liệu & Huang Luotian \\
37 & Bàn về nhân tử Lagrange và đối ngẫu trong OI & Shi Kaicheng \\
54 & Thuật toán liên quan tới bài toán cập nhật đoạn và truy vấn đoạn cùng ứng dụng & Shen Jiyu \\
69 & Bàn về một số phương pháp giải bài toán quy hoạch động trong OI & Cao Li \\
85 & Bàn về một số cách xử lý hợp nhất thông tin & Feng Zhengwei \\
100 & Một số phương pháp cho bài toán đếm thứ tự topo trên cây & Chen Xinyang \\
121 & Bàn về tối ưu hằng số trên phần cứng hiện đại & Song Jiaxing \\
171 & Bàn về ứng dụng của một lớp phương pháp bao hàm--loại trừ và nghịch đảo & Dong Yiqi \\
182 & Một lớp kỹ thuật dùng hàm sinh để mô tả thời điểm dừng của quá trình ngẫu nhiên & Ye Kai \\
196 & Bàn về bài toán đồ thị đường trong thi tin học & Zhai Jincheng \\
206 & Bàn về một lớp bài toán đoạn trên cây & Zhu Yifan \\
217 & Báo cáo ra đề \emph{Thế giới chìm trong cổ tích ngủ say} & Ye Liqi \\
230 & Bài toán matching trên đồ thị có trọng số đỉnh & Ke Yisi \\
253 & Bàn lại về đếm đường đi trên lưới & Wu Chang \\
277 & Ứng dụng liên phân số và thuật toán Euclid trong OI & Ke Yijing \\
299 & Bàn về độ phức tạp và ứng dụng của nó trong giải bài toán & Yang Minxing \\
315 & Bàn về thuật toán duy trì một loại cặp số nguyên tố cùng nhau và ước chung lớn nhất & Du Guancheng \\
321 & Bàn về lý thuyết tính toán và các bài toán khó trong OI & Yan Chenxiao \\
336 & Bàn về ứng dụng thuật toán lattice trong phân tích đa thức hệ số nguyên & Xie Zihan \\
358 & Bàn về bài toán tô màu danh sách của đồ thị & Zhong Zihou \\
367 & Từ \emph{Kén}: phân tích phương pháp quan sát trong thi tin học qua ví dụ & Li Jiayou \\
376 & Bàn về \(q\)-analogue & Li Mingleyang \\
400 & Bàn về chu trình tỉ lệ nhỏ nhất và ứng dụng & Wei Jiaze
\end{tabular}
\end{center}

\section*{Bài 1. Một số tiến triển về bài toán tính tổng hàm nhân tính}
\addcontentsline{toc}{section}{Bài 1. Một số tiến triển về bài toán tính tổng hàm nhân tính}

\begin{center}
Zhou Kangyang, Trường Văn Uyên thuộc Tập đoàn Giáo dục Trung học Học Quân,
Hàng Châu
\end{center}

\subsection*{Tóm tắt}

Tính tổng hàm nhân tính là một lớp bài toán quan trọng trong số học. Với bài
toán này, các nghiên cứu trước đây chủ yếu hướng tới giảm số lượng thừa số
\(\log\) trong độ phức tạp \(\widetilde O(n^{2/3})\).\footnote{Một nghiên cứu
liên quan: \url{https://negiizhao.blog.uoj.ac/blog/7165}.} Sau khi suy nghĩ,
tác giả nhận thấy có thể đạt độ phức tạp thời gian
\(O(n^{1/2}\operatorname{poly}\log n)\).\footnote{Bản đầu được viết ngày
2023-07-11 và công bố tại
\url{https://www.cnblogs.com/zkyJuruo/p/17544928.html}.}

Bài viết này giới thiệu thuật toán đó và các tối ưu của nó.

\subsection*{1. Kiến thức chuẩn bị}

\subsubsection*{1.1. Định nghĩa và quy ước}

\paragraph{\term{Định nghĩa 1.1}.}
Hàm số học là hàm có miền xác định là tập số nguyên dương và miền giá trị là
tập số phức.

Với một hàm số học \(f\), ta quy ước
\[
  S_f(n)=\sum_{i=1}^{n} f(i).
\]

\paragraph{\term{Định nghĩa 1.2}.}
Hàm nhân tính là một hàm số học \(f(n)\) xác định trên các số nguyên dương,
thỏa
\[
  f(1)=1,\qquad
  \gcd(a,b)=1 \Longrightarrow f(ab)=f(a)f(b).
\]

Với một hàm nhân tính \(f\), chỉ cần biết mọi giá trị \(f(p^k)\) với
\(p\in\mathbb P\), \(k\in\mathbb Z^+\), ta có thể xác định toàn bộ hàm \(f\).
Một số hàm nhân tính thường gặp có giá trị trên \(p^k\) như sau:
\begin{itemize}
  \item hàm \(\epsilon\) (hàm đơn vị), thỏa \(\epsilon(p^k)=0\);
  \item hàm \(\mu\), thỏa \(\mu(p^k)=-[k=1]\);
  \item hàm \(\varphi\), thỏa \(\varphi(p^k)=(p-1)p^{k-1}\);
  \item hàm \(I\), thỏa \(I(p^k)=1\);
  \item hàm \(\operatorname{id}_t\), thỏa
  \(\operatorname{id}_t(p^k)=p^{kt}\).
\end{itemize}

\paragraph{\term{Định nghĩa 1.3}.}
Với hai hàm số học \(f,g\), định nghĩa tích chập Dirichlet của chúng là
\[
  (f*g)(x)=\sum_{d\mid x} f(d)g\left(\frac{x}{d}\right).
\]

Nếu xem tích chập Dirichlet là phép nhân của các hàm số học, và phép cộng trực
tiếp của hàm là phép cộng, thì phép cộng và phép nhân trên các hàm số học thỏa
giao hoán, kết hợp và phân phối.

\paragraph{\term{Định lý 1.1}.}
Tích chập Dirichlet \(h\) của hai hàm nhân tính \(f,g\) cũng là một hàm nhân
tính.

\paragraph{\term{Chứng minh}.}
Chỉ cần chứng minh với mọi \(\gcd(x,y)=1\) ta có \(h(xy)=h(x)h(y)\):
\[
\begin{aligned}
  h(xy)
  &= \sum_{d\mid xy} f(d)g\left(\frac{xy}{d}\right) \\
  &= \sum_{u\mid x,\ v\mid y} f(uv)g\left(\frac{xy}{uv}\right) \\
  &= \sum_{u\mid x} f(u)g\left(\frac{x}{u}\right)
     \sum_{v\mid y} f(v)g\left(\frac{y}{v}\right) \\
  &= h(x)h(y).
\end{aligned}
\]

\subsubsection*{1.2. Powerful Number}

\paragraph{\term{Định nghĩa 1.4}.}
Với một số nguyên dương \(m\), nếu với mọi thừa số nguyên tố \(p\) của \(m\) đều
có \(p^2\mid m\), thì gọi \(m\) là một \emph{Powerful Number}.

Trong phần sau, ta viết tắt \emph{Powerful Number} là PN. Trong công thức,
\(PN\) biểu thị tập tất cả các PN.

\paragraph{\term{Định lý 1.2}.}
Mọi PN đều có thể biểu diễn dưới dạng \(a^2b^3\), trong đó \(a,b\) là các số
nguyên dương.

\paragraph{\term{Chứng minh}.}
Giả sử
\[
  x=\prod_i p_i^{\alpha_i}
\]
là một PN. Chỉ cần biểu diễn từng \(p_i^{\alpha_i}\) dưới dạng
\(a_i^2b_i^3\), rồi lấy \(a=\prod_i a_i\), \(b=\prod_i b_i\).

Nếu \(\alpha_i\) lẻ, lấy
\[
  a_i=p_i^{(\alpha_i-3)/2},\qquad b_i=p_i;
\]
nếu \(\alpha_i\) chẵn, lấy
\[
  a_i=p_i^{\alpha_i/2},\qquad b_i=1.
\]

\paragraph{\term{Định lý 1.3}.}
Số lượng PN không vượt quá \(n\) là \(O(\sqrt n)\).

\paragraph{\term{Chứng minh}.}
Biểu diễn mọi PN không vượt quá \(n\) dưới dạng \(a^2b^3\). Khi liệt kê \(a\),
số PN không vượt quá
\[
  \sum_{a=1}^{\sqrt n}\left(\frac{n}{a^2}\right)^{1/3}.
\]
Tổng này có thể ước lượng bởi
\[
  \int_1^{\sqrt n}\left(\frac{n}{x^2}\right)^{1/3}\,dx
  = n^{1/3}\int_1^{\sqrt n} x^{-2/3}\,dx
  = n^{1/3}(3n^{1/6}-3),
\]
tức ở cấp \(O(\sqrt n)\).

Sàng tuyến tính mọi số nguyên tố không vượt quá \(\sqrt n\), rồi DFS trên các
số nguyên tố sẽ tìm được mọi PN không vượt quá \(n\). Độ phức tạp bằng số PN,
tức \(O(\sqrt n)\).

\paragraph{\term{Ví dụ 1.1} (phần \(n=1\) của UOJ Long Round 7,
\emph{Mã kiểm tra}).}
Cho \(c\). Hàm nhân tính \(q(x)\) thỏa
\[
  q(p^k)=p^{c\lfloor k/2\rfloor}.
\]
Tính \(\sum_{i=1}^{m} q(i)\), đáp án lấy modulo \(2^{32}\), với
\(m\le 1.2\times 10^{11}\).

Thiết kế hàm nhân tính \(f(x)\) thỏa \(f(p)=q(p)\), và dựng hàm nhân tính
\(g\) thỏa \(g*f=q\). Với số nguyên tố \(p\),
\[
  g(p)f(1)+g(1)f(p)=q(p),
\]
suy ra \(g(p)=0\). Vì vậy \(g\) chỉ có giá trị khác không tại các PN.

Ta có
\[
\begin{aligned}
  \sum_{i=1}^{m} q(i)
  &= \sum_{i=1}^{m}\sum_{jk=i} g(j)f(k) \\
  &= \sum_{\substack{j\in PN\\j\le m}} g(j)
     \sum_{k=1}^{\lfloor m/j\rfloor} f(k) \\
  &= \sum_{\substack{j\in PN\\j\le m}} g(j)\left\lfloor\frac{m}{j}\right\rfloor.
\end{aligned}
\]

Tiền xử lý \(g(p^k)\) cho từng \(p^k\). Từ
\[
  q(p^k)=\sum_{i=0}^{k} f(p^i)g(p^{k-i})
\]
có thể giải \(g(p^k)\) theo thứ tự tăng dần của \(k\), đồng thời duy trì giá
trị hàm \(g\) trong lúc DFS tìm PN. Như vậy bài toán được giải trong
\(O(\sqrt m)\).

\subsubsection*{1.3. Sàng Du Jiao}

Với một số hàm nhân tính đặc biệt như \(\mu\) và \(\varphi\), ta có thể xây
dựng tích chập Dirichlet để tính tổng tiền tố trong thời gian \(O(n^{2/3})\).
Trong cộng đồng thi tin học Trung Quốc, thuật toán kiểu này được gọi là
``sàng Du Jiao''.

Với hàm \(\mu\), ta có \(\mu*I=\epsilon\).

Xét bài toán tổng quát hơn. Nếu có các hàm số học \(A,B,C\) thỏa
\[
  A(1)=B(1)=C(1)=1,\qquad A*B=C,
\]
và ta tính nhanh được các giá trị điểm của \(S_B\), \(S_C\), thì làm sao tính
\(S_A(n)\)?

Chú ý rằng
\[
\begin{aligned}
  \sum_{i=1}^{n} C(i)
  &= \sum_{i=1}^{n}\sum_{jk=i} A(j)B(k) \\
  &= \sum_{k=1}^{n} B(k)\sum_{j=1}^{\lfloor n/k\rfloor} A(j) \\
  &= S_A(n)+\sum_{k=2}^{n} B(k)S_A\left(\left\lfloor\frac{n}{k}\right\rfloor\right).
\end{aligned}
\]
Do đó
\[
  S_A(n)=S_C(n)-\sum_{k=2}^{n}
  B(k)S_A\left(\left\lfloor\frac{n}{k}\right\rfloor\right),
\]
có thể tính \(S_A(n)\) bằng đệ quy. Lại có
\[
  \left\lfloor\frac{\lfloor n/a\rfloor}{b}\right\rfloor
  = \left\lfloor\frac{n}{ab}\right\rfloor,
\]
nên trong quá trình đệ quy, các giá trị thật sự cần tính chỉ là
\[
  S_A\left(\left\lfloor\frac{n}{k}\right\rfloor\right)
  \qquad (k\in\mathbb Z^+).
\]

\paragraph{\term{Định lý 1.4}.}
Với số nguyên dương \(n\), tập
\[
  \left\{\left\lfloor\frac{n}{i}\right\rfloor\mid i\in\mathbb Z^+\right\}
\]
có kích thước \(O(\sqrt n)\).

\paragraph{\term{Chứng minh}.}
Với \(k\le\sqrt n\), chỉ có \(\sqrt n\) giá trị
\(\lfloor n/k\rfloor\), hiển nhiên là \(O(\sqrt n)\) loại. Với
\(k>\sqrt n\), ta có \(\lfloor n/k\rfloor\le\sqrt n\), nên cũng chỉ có
\(O(\sqrt n)\) loại.

Vì vậy ta chỉ cần tính \(O(\sqrt n)\) giá trị điểm của \(S_A\).

Khi tính một \(S_A(m)\), gom các đoạn liên tiếp có cùng giá trị
\(\lfloor m/k\rfloor\) lại với nhau (tức chia khối theo phép chia nguyên), ta
có thể đệ quy tới các bài toán con trong \(O(\sqrt m)\).

Với \(m\le n^{2/3}\), ta xử lý trước \(n^{2/3}\) giá trị điểm đầu của \(A\).
Với \(m>n^{2/3}\), dùng công thức đệ quy ở trên. Phần này có độ phức tạp
\[
  O\left(\sum_{i=1}^{n^{1/3}}\sqrt{\frac{n}{i}}\right)
  = O\left(\int_1^{n^{1/3}}\sqrt{\frac{n}{x}}\,dx\right)
  = O(n^{2/3}).
\]
Tổng độ phức tạp là \(O(n^{2/3})\) cộng với phần xử lý trước \(O(n^{2/3})\) giá
trị đầu của \(A\). Với \(\mu\) và \(\varphi\), phần xử lý trước này đều làm
được trong \(O(n^{2/3})\).

\paragraph{\term{Định nghĩa 1.5}.}
Định nghĩa \term{sàng khối} của hàm số học \(f\) đối với \(n\) là các giá trị
điểm của \(S_f(n)\) trên tập
\[
  \left\{\left\lfloor\frac{n}{i}\right\rfloor\mid i\in\mathbb Z^+\right\}.
\]

Trong quá trình trên, ta dùng sàng khối của \(B\) và \(C\) đối với \(n\) để
tìm sàng khối của \(A\) đối với \(n\). Hiển nhiên, nếu biết sàng khối của
\(A\) và \(B\) đối với \(n\), ta cũng có thể tìm sàng khối của \(C\) đối với
\(n\). Ta gọi \(C\) là \term{tích chập sàng khối} của \(A\) và \(B\). Các phép
toán trên sàng khối vẫn thỏa giao hoán, kết hợp và phân phối.

\subsection*{2. Tích chập sàng khối}

Với các hàm số học \(A,B,C\) thỏa
\[
  A(1)=B(1)=C(1)=1,\qquad A*B=C,
\]
sàng Du Jiao ở trên đã dùng sàng khối của \(B,C\) đối với \(n\) để giải sàng
khối của \(A\) đối với \(n\).

Ở đây trước hết xét bài toán đơn giản hơn: đã biết sàng khối của \(A,B\) đối
với \(n\), hãy tính sàng khối của \(C=A*B\) đối với \(n\). Bài toán này thực
ra có thể giải trong \(O(n^{1/2}\operatorname{poly}\log n)\). Trong phần dưới,
mặc định sàng khối là sàng khối đối với \(n\).

Ta cần tính các tổng tiền tố trên mọi \(\lfloor n/t\rfloor\) của
\[
  C(z)=\sum_{xy=z} A(x)B(y).
\]

Trước hết xét trường hợp \(x>\sqrt n\); trường hợp \(y>\sqrt n\) là đối xứng.
Nếu \(xy\le\lfloor n/t\rfloor\) thì \(x\le n/(ty)\). Do đó khi liệt kê
\(t,y\), các \(x\) có đóng góp là một đoạn tiền tố. Phần này có độ phức tạp
\(O(\sqrt n\log n)\).

Như vậy chỉ còn cần xử lý trường hợp \(x\le\sqrt n\) và \(y\le\sqrt n\).
Điều kiện \(xy\le n/t\) tương đương
\[
  \ln x+\ln y\le \ln\left(\frac{n}{t}\right).
\]
Ta làm một ước lượng như sau. Chọn một độ dài khối nguyên \(S\) với
\(3\le S\le n\), và cho rằng
\[
  \ln x+\ln y\le \ln\left(\frac{n}{t}\right)
\]
khi và chỉ khi
\[
  \lfloor S\ln x\rfloor+\lfloor S\ln y\rfloor
  \le S\ln\left(\frac{n}{t}\right).
\]
Ước lượng này có thể làm bằng nhân đa thức. Thiết kế hai đa thức \(F(z),G(z)\)
sao cho \([z^k]F(z)\) là tổng các \(A(x)\) thỏa
\(\lfloor S\ln x\rfloor=k\), và \([z^k]G(z)\) là tổng các \(B(x)\) thỏa
\(\lfloor S\ln x\rfloor=k\). Tính \(H(z)=F(z)G(z)\), khi đó hệ số
\([z^k]H(z)\) là tổng \(A(x)B(y)\) với
\(\lfloor S\ln x\rfloor+\lfloor S\ln y\rfloor=k\). Chỉ cần làm một phép chập
độ dài \(S\ln n\).

Nhưng làm như vậy hiển nhiên không đúng hoàn toàn. Nó chỉ sai khi
\[
  \ln x+\ln y\le \ln\left(\frac{n}{t}\right)
  \quad\text{và}\quad
  \lfloor S\ln x\rfloor+\lfloor S\ln y\rfloor
  > S\ln\left(\frac{n}{t}\right).
\]
Khi đó
\[
  S\ln x+S\ln y\in
  \left(S\ln\left(\frac{n}{t}\right)-2,\,
        S\ln\left(\frac{n}{t}\right)\right],
\]
nên
\[
  \ln x+\ln y+\ln t\in
  \left(\ln n-\frac{2}{S},\,\ln n\right],
\]
tức
\[
  xyt\in (ne^{-2/S},n].
\]
Trong đó \(e^{-2/S}\) ở cấp \(1-O(1/S)\). Vì vậy \(xyt\) nằm trong một đoạn có
độ dài \(O(n/S)\).

Giả sử đoạn này là \([n-L,n]\). Dùng sàng đoạn để lấy phân tích thừa số nguyên
tố của mọi số trong đoạn. Cụ thể hơn, trước hết sàng mọi số nguyên tố nhỏ không
vượt quá \(\sqrt n\). Với mỗi số nguyên tố, ta tìm được các bội của nó trong
\([n-L,n]\), từ đó sàng ra mọi thừa số nguyên tố \(\le\sqrt n\) của các số
trong đoạn. Một số có nhiều nhất một thừa số nguyên tố \(>\sqrt n\), nên chỉ
cần xét phần còn lại sau khi chia hết các số nguyên tố nhỏ.

Sau khi có phân tích thừa số nguyên tố, DFS trên các thừa số nguyên tố sẽ nhanh
chóng tìm được mọi bộ \((x,y,t)\) có thể. Với một bộ \((x,y,t)\), kiểm tra trong
\(O(1)\) xem nó có bị ước lượng sai hay không. Với một
\(v\in[n-L,n]\), thời gian đóng góp là \(d_3(v)\), trong đó
\[
  d_3(v)=\sum_{xyz=v}1.
\]
Theo kết quả trong số học giải tích \([1]\),
\[
  \sum_{i\le n} d_3(i)=nP_3(\ln n)+O(n^{43/96+\epsilon}),
\]
nên
\[
  \sum_{i=n-L}^{n} d_3(i)=LP_3(\ln n)+O(n^{43/96+\epsilon}),
\]
trong đó \(P_3\) là một đa thức bậc hai. Bản thân đầu ra của thuật toán đã có
\(O(\sqrt n)\), lớn hơn \(O(n^{43/96+\epsilon})\), nên có thể bỏ qua phần độ
phức tạp này. Phần còn lại có độ phức tạp
\[
  O(L\log^2 n)=O\left(\frac{n}{S}\log^2 n\right).
\]
Tổng độ phức tạp là
\[
  O\left(S\log^2 n+\frac{n}{S}\log^2 n\right).
\]
Lấy \(S=\sqrt n\) thu được độ phức tạp thời gian
\[
  O(\sqrt n\log^2 n).
\]

\subsection*{3. Tính tổng hàm nhân tính}

Với một hàm nhân tính \(f\) thỏa một số tính chất đặc biệt, chẳng hạn
\[
  f(p^k)=p^k(p^k-1),
\]
làm sao tính sàng khối của \(f\)?

Trước hết, giá trị \(f(p^k)\) với \(k\ge 2\) không quá quan trọng. Nếu có một
hàm nhân tính khác \(g\) thỏa \(f(p)=g(p)\), và tổng tiền tố của \(g\) dễ tính,
thì chỉ cần thiết kế hàm nhân tính \(h\) thỏa \(h*g=f\). Theo kết luận ở trên,
\(h\) chỉ có giá trị tại PN, và các giá trị đó có thể tính trong \(O(\sqrt n)\).
Vì vậy dễ dàng tính sàng khối của \(h\) trong \(O(\sqrt n)\), rồi dùng tích
chập sàng khối của \(h\) và \(g\) để tính sàng khối của \(f\).

Ta trước hết thiết kế một hàm số học \(q\) chỉ khác không tại các số nguyên tố
và thỏa \(q(p)=f(p)\).

Với một hàm số học \(q\) có \(q(1)=0\), định nghĩa
\[
  \exp(q)=\sum_{i=0}^{\infty}\frac{q^i}{i!},
\]
trong đó phép chập là tích chập sàng khối. Xét \(d=\exp(q)\). Với
\[
  x=\prod_{i=1}^{k} p_i^{\alpha_i},
\]
hàm \(d(x)\) chỉ có đóng góp từ hạng
\[
  q^{\sum_{i=1}^{k}\alpha_i},
\]
và giá trị bằng
\[
  \frac{\prod_{i=1}^{k} q(p_i)^{\alpha_i}}
       {\left(\sum_{i=1}^{k}\alpha_i\right)!}
  \binom{\sum_{i=1}^{k}\alpha_i}{\alpha_1,\alpha_2,\ldots,\alpha_k}
  =
  \frac{\prod_{i=1}^{k} q(p_i)^{\alpha_i}}
       {\prod_{i=1}^{k}\alpha_i!}.
\]

Do đó \(d\) là một hàm nhân tính và
\[
  d(p^k)=\frac{f(p)^k}{k!},
\]
đặc biệt \(d(p)=f(p)\). Như vậy, nếu tính được sàng khối của \(q\), ta có thể
dùng \(\log n\) lần tích chập sàng khối (khi \(i>\log n\), \(q^i\) không còn
giá trị trong \(n\) hạng đầu) để tính sàng khối của \(d\), từ đó suy ra sàng
khối của \(f\).

Làm sao tính sàng khối của \(q\)? Xét một ý tưởng tương tự sàng Min\_25. Tách
hàm \(q\) thành tổng có trọng số của vài \(q_i\), rồi dùng sàng khối của các
\(q_i\) để tính sàng khối của \(q\). Ví dụ, nếu
\[
  q(p)=p(p-1),
\]
ta có thể tách thành
\[
  q_1(p)=p^2,\qquad q_2(p)=p,\qquad q=q_1-q_2.
\]

Khi \(q(p)\) là đa thức bậc hằng theo \(p\), ta đều có thể tách như vậy.

Làm sao tính sàng khối của \(q_i\)? Ta làm ngược lại. Lấy \(q_1(p)=p^2\) làm
ví dụ. Với hàm nhân tính \(d\) thỏa
\[
  d(p^k)=q_1(p)^k,
\]
ta có \(d(x)=x^2\), nên sàng khối rất dễ tính. Với hàm nhân tính \(d\) thỏa
\[
  d(p^k)=\frac{q_1(p)^k}{k!},
\]
do nó có cùng giá trị tại các điểm \(p\) với hàm trước, ta có thể dùng PN để
chuyển đổi lẫn nhau trong độ phức tạp của một lần tích chập sàng khối. Vì vậy
sàng khối của nó cũng tính được.

Hàm \(d\) thỏa \(d(p^k)=q_1(p)^k/k!\) chính là \(\exp(q_1)\). Điều này gợi ý
dùng một ý tưởng tương tự \(\ln\) để suy ngược \(q_1\) từ \(d\).

Với một hàm số học \(d\) có \(d(1)=1\), định nghĩa
\[
  \ln(d)=\sum_{i=1}^{\infty}
  \frac{(-1)^{i-1}(d-\epsilon)^i}{i}.
\]
Ta có \(q_1=\ln(d)\).

\paragraph{\term{Chứng minh}.}
Chỉ cần lần lượt chứng minh \(\exp(\ln(d))=d\) và hàm \(q\) thỏa
\(\exp(q)=d\) là duy nhất.

Trước hết chứng minh \(\exp(\ln(d))=d\):
\[
  \exp(\ln(d))
  =
  \sum_{j=0}^{\infty}
  \frac{\left(\sum_{i=1}^{\infty}
  \frac{(-1)^{i-1}(d-\epsilon)^i}{i}\right)^j}{j!}.
\]
Dùng tính phân phối và kết hợp của sàng khối để khai triển, ta sẽ nhận được
biểu thức dạng \(\sum_i a_i d^i\). Có thể dùng đa thức \(A(z)\) để mô tả nó,
trong đó \([z^k]A(z)=a_k\). Định nghĩa của \(\exp\) và \(\ln\) trong tích chập
sàng khối tương ứng với đa thức:
\[
  \exp(z)=\sum_{i=0}^{\infty}\frac{z^i}{i!},\qquad
  \ln(1+z)=\sum_{i=1}^{\infty}\frac{(-1)^{i-1}z^i}{i}.
\]
Vì \(A(z)=\exp(\ln z)=z\), nên \(a_i=[i=1]\), do đó \(\exp(\ln(d))=d\).

Nếu \(\exp(q)=d\), tức
\[
  \epsilon+q+\sum_{i=2}^{\infty}\frac{q^i}{i!}=d,
\]
ta có thể xác định lần lượt giá trị \(q(i)\) theo thứ tự tăng dần. Đặt
\[
  p=\sum_{i=2}^{\infty}\frac{q^i}{i!},
\]
khi đó \(p(i)\) chỉ liên quan tới \(q(j)\) với \(j<i\), nên có thể dùng
\[
  q(i)=d(i)-[i=1]-p(i)
\]
để giải \(q(i)\). Do đó \(q\) được xác định duy nhất.

Trong công thức \(\ln\) ở trên cũng chỉ cần liệt kê \(i\le\log n\), vì vậy
cũng chỉ cần \(\log n\) lần tích chập sàng khối.

Do đó thuật toán này có độ phức tạp thời gian
\[
  O(\sqrt n\log^3 n).
\]

\subsection*{4. Tối ưu}

Dù ta đã có một thuật toán \(O(n^{1/2}\operatorname{poly}\log n)\), nó vẫn quá
chậm để dùng thực tế trong OI. Có tối ưu được không?

Dưới đây là một số hướng tối ưu đơn giản, có thể hạ độ phức tạp thời gian của
thuật toán xuống
\[
  O\left(\frac{\sqrt n\log^2 n}{\sqrt{\log\log n}}\right).
\]

\subsubsection*{4.1. Tối ưu \(\exp\)}

Khi tính \(\exp(f)\), với \(f\) chỉ có giá trị tại các số nguyên tố, ta phải
làm \(\log n\) lần tích chập sàng khối; điều này khá lãng phí. Xét cách tối
ưu.

Trước hết tách \(f\) thành phần có chỉ số \(\le\sqrt n\), gọi là \(f_0\), và
phần có chỉ số \(>\sqrt n\), gọi là \(f_1\). Khi đó
\[
  \exp(f)=\exp(f_0+f_1)=\exp(f_0)\exp(f_1).
\]
Vì \(\exp(f_1)=\epsilon+f_1\) dễ tính, phần khó là \(\exp(f_0)\).

Khi tính \(\exp(f_0)\), \( \sqrt n \) hạng đầu của kết quả cuối có thể tính
trực tiếp; khó khăn nằm ở phần \(>\sqrt n\).

Xét cách tính trực tiếp bằng phương pháp lấy phần nguyên rồi sửa sai ở trên.
Chọn số nguyên \(S\), và cộng \(f_{0,i}\) vào vị trí
\(v_{\lfloor S\ln i\rfloor}\). Sau đó trực tiếp lấy \(\exp\) đa thức theo
\(v\).

Khi nào cách này sai? Giả sử
\[
  m=\prod_{i=1}^{k}p_i^{\alpha_i}.
\]
Nếu vị trí của \(m\) bị ước lượng sai, đặt \(t=\lfloor n/m\rfloor\), thì
\[
  \sum_{i=1}^{k}\alpha_i\lfloor S\ln p_i\rfloor
  \in
  \left(S\ln\left(\frac{n}{t}\right),\,
        S\ln\left(\frac{n}{t}\right)+\sum_{i=1}^{k}\alpha_i\right).
\]
Suy ra
\[
  \sum_{i=1}^{k}\alpha_i\ln p_i+\ln t-\ln n
  \in
  \left[-\frac{\sum_{i=1}^{k}\alpha_i}{S},\,0\right],
\]
tức
\[
  \ln(mt)\in
  \left[\ln n-\frac{\sum_i\alpha_i}{S},\,\ln n\right].
\]
Đây là một đoạn có độ dài
\[
  O\left(\frac{n\log n}{S}\right),
\]
vì vậy chỉ các \(mt\) trong đoạn này mới bị ước lượng sai.

Dùng sàng đoạn để sàng mọi thừa số nguyên tố \(\le\sqrt n\) trong đoạn đó. Sau
khi sàng ra thừa số nguyên tố, DFS trên các thừa số nguyên tố sẽ nhanh chóng
tìm các cặp \((m,t)\); trong lúc DFS cũng có thể đồng thời tìm giá trị hàm tại
\(m\). Độ phức tạp DFS cùng cấp với số cặp.

Số cặp không vượt quá
\[
  O\left(\sum_{m=1}^{\sqrt n}\frac{n\log n}{Sm}\right)
  =
  O\left(\frac{n\log^2 n}{S}\right).
\]
Phần đa thức phía trước có độ phức tạp \(O(S\log^2 n)\). Lấy \(S=\sqrt n\) thu
được độ phức tạp \(O(\sqrt n\log^2 n)\).

Ta còn có thể chọn một ngưỡng \(L\) và xử lý riêng các thừa số nguyên tố
\(\le L\). Khi tính \(\exp(f_0)\), không xét các thừa số \(\le L\), tức đặt
chúng về \(0\) ở các vị trí tương ứng trong \(f_0\). Sau khi tính xong, thêm
lần lượt các thừa số nguyên tố đó vào.

Khi thêm số nguyên tố \(p\), giả sử đáp án trước đó là \(d\), thì đáp án mới
là
\[
  d'(n)=\sum_{k=0}^{\log_p n}
  f(p^k)d\left(\left\lfloor\frac{n}{p^k}\right\rfloor\right).
\]
Chi phí cần thiết là \(\sqrt n\log_p n\). Tổng chi phí thêm vào là
\[
  \sum_{\substack{p\le L\\p\in\mathbb P}}
  \frac{\sqrt n\log n}{\log p}
  =
  O\left(\frac{L\sqrt n\log n}{\log L}\right).
\]

Sau khi loại bỏ các thừa số nguyên tố \(\le L\), mỗi \(m\) ở trên chỉ còn thừa
số nguyên tố \(>L\). Nếu
\[
  m=\prod_{i=1}^{k}p_i^{\alpha_i},
\]
thì
\[
  \sum_{i=1}^{k}\alpha_i\le\log_L n.
\]
Vì vậy độ dài đoạn có thể ước lượng sai giảm xuống
\[
  O\left(\frac{n\log n}{S\log L}\right).
\]
Khi đó độ phức tạp tính \(\exp(f_0)\) trở thành
\[
  O\left(\frac{n\log n\log_L n}{S}+S\log^2 n\right).
\]
Lấy
\[
  S=\sqrt{\frac{n}{\log L}},
\]
ta được độ phức tạp
\[
  O\left(\frac{\sqrt n\log^2 n}{\sqrt{\log L}}\right).
\]
Lấy \(L=\log n\), độ phức tạp thời gian là
\[
  O\left(\frac{\sqrt n\log^2 n}{\sqrt{\log\log n}}\right).
\]

\subsubsection*{4.2. Tối ưu \(\ln\)}

Khi tính \(\ln\), ta cần biến một hàm nhân tính thỏa
\[
  f(p^k)=\frac{f(p)^k}{k!}
\]
thành một hàm số học \(g\) chỉ có giá trị tại \(p\) và thỏa \(g(p)=f(p)\).

Từ tính chất ở trên có thể trực tiếp lấy \(\sqrt n\) hạng đầu của
\(g=\ln(f)\). Gọi hàm số học chỉ giữ \(\sqrt n\) hạng đầu của \(g\) là \(g_0\),
và đặt \(g_1=g-g_0\). Khi đó
\[
  \exp(g_0+g_1)=\exp(g_0)\exp(g_1).
\]
Chú ý rằng tích chập của hai \(g_1\) có hạng thấp nhất \(>n\), nên mọi giá trị
trên sàng khối đều bằng \(0\) và không cần tính. Vì vậy
\[
  f=\exp(g_0)+\exp(g_0)g_1.
\]
Sau khi tính \(h=\exp(g_0)\), chỉ cần giải
\[
  f-h=hg_1.
\]

Bây giờ ta cần giải một bài toán chia trên sàng khối. Với các hàm số học
\(A,B,C\) thỏa \(B(1)=1\) và \(A*B=C\), làm sao lấy sàng khối của \(A\) từ
sàng khối của \(B,C\)?

Trước hết, \(\sqrt n\) hạng đầu của \(A\) dễ giải trong
\(O(\sqrt n\log n)\). Đặt \(A=A_0+A_1\), trong đó \(A_0\) chỉ giữ giá trị của
\(A\) tại các vị trí \(\le\sqrt n\), còn \(A_1\) chỉ giữ giá trị tại các vị trí
\(>\sqrt n\). Khi đó
\[
  C=A_0B+A_1B.
\]
Dùng một lần tích chập sàng khối để lấy \(A_0B\), đặt \(D=C-A_0B\); ta chỉ cần
giải \(A_1B=D\).

Với mọi \(k\),
\[
  D\left(\left\lfloor\frac{n}{k}\right\rfloor\right)
  =
  \sum_{i=1}^{\infty} B(i)
  S_{A_1}\left(\left\lfloor\frac{n}{ik}\right\rfloor\right).
\]
Vì \(S_{A_1}(\lfloor n/(ik)\rfloor)\) chỉ có giá trị khi
\(ik\le\sqrt n\), tổng số \(i\) cần liệt kê là
\[
  \sum_{k=1}^{\sqrt n}\left\lfloor\frac{\sqrt n}{k}\right\rfloor
  =
  O(\sqrt n\log n).
\]
Do đó thời gian là \(O(\sqrt n\log n)\).

Tới đây, ta đã thực hiện được phép \(\ln\) trên sàng khối bằng một lần
\(\exp\) của hàm chỉ có giá trị tại số nguyên tố, một lần tích chập sàng khối,
và thêm \(O(\sqrt n\log n)\) độ phức tạp phụ.

\subsubsection*{4.3. Tối ưu tích chập sàng khối}

Lúc này nút thắt độ phức tạp lại nằm ở tích chập sàng khối. Tích chập sàng
khối có thể làm tốt hơn không? Cũng có thể.

Trong phần sửa sai của ước lượng tích chập sàng khối, với mỗi
\(m\in[n-L,n]\), ta cần liệt kê mọi bộ ba \((x,y,t)\) thỏa \(xyt=m\). Cố định
\(m\), liệu phần này có thể tốt hơn \(d_3(m)\)?

Quan sát các giá trị ta thật sự quan tâm:
\begin{itemize}
  \item giá trị của \(xy\) và \(t\), vì chúng ảnh hưởng tới vị trí cần sửa;
  \item giá trị của \(\lfloor S\ln x\rfloor+\lfloor S\ln y\rfloor\), vì nó
  ảnh hưởng tới việc vị trí \(xy\) có được cập nhật hay không.
\end{itemize}

Chú ý tính chất
\[
  S\ln(xy)\le
  \lfloor S\ln x\rfloor+\lfloor S\ln y\rfloor
  < S\ln(xy)+2.
\]
Vì vậy khi đã xác định \(xy\), giá trị
\(\lfloor S\ln x\rfloor+\lfloor S\ln y\rfloor\) chỉ có hai khả năng. Do đó với
\(x,y\), ta chỉ cần biết
\[
  w_x=\lfloor S\ln x\rfloor\bmod 2,\qquad
  w_y=\lfloor S\ln y\rfloor\bmod 2.
\]

Ta cần làm một phép chập để các đóng góp \(x,y\) đi vào \(xy\). Giả sử phân
tích thừa số nguyên tố của \(m\) là
\[
  m=\prod_{i=1}^{k}p_i^{\alpha_i}.
\]
Phép chập ở trên thực chất là một phép chập \(k\)-chiều; miền giá trị của chiều
thứ \(i\) là \([0,\alpha_i]\). Tuy phép chập còn chịu ảnh hưởng của \(w\), chỉ
cần liệt kê bốn trường hợp \(w_x=0/1\), \(w_y=0/1\) và làm phép chập nhiều
chiều thông thường, hoặc dùng phản diễn đơn vị căn bậc hai: với
\(c=1,-1\), dùng
\[
  c^{w_x}c^{w_y}=c^{(w_x+w_y)\bmod 2}
\]
để giảm còn hai phép chập.

Bài toán chập đa thức nhiều biến có thể tham khảo luận văn đội tuyển tập huấn
năm 2021 của Li Baitian \([2]\). Thời gian một lần là
\[
  d(m)\omega(m)\log d(m).
\]

Xét cách phân tích độ phức tạp. Chọn một hằng số \(l\). Số thừa số nguyên tố
\(>n^{1/l}\) chỉ là \(O(l)\), và đóng góp của chúng vào \(\omega(m)\) cũng như
\(\log d(m)\) đều không quá hằng số, nên có thể bỏ qua.

Với hạng \(\log d(m)\), có thể dùng
\[
  \sum_{i=1}^{k}\sum_{e=1}^{\alpha_i}\frac1e
\]
để ước lượng. Ngoài ra, nếu \(p^t\Vert m\), thì
\[
  (t+1)d\left(\frac{m}{p^t}\right)\le d(m),
\]
nên có thể ước lượng \(d(ipq^e)\) với \(p,q\in\mathbb P\) bởi \(O(d(i)e)\).

Với một \(m\), phần \(d(m)\log d(m)\omega(m)\) có thể viết ở cấp
\[
  O\left(d(m)
  \left(\sum_{\substack{p\mid m\\p\in\mathbb P}}1\right)
  \left(\sum_{\substack{q\in\mathbb P,\ e\ge 1\\q^e\Vert m}}\frac1e\right)
  \right).
\]
Trường hợp \(p=q\) có thể bỏ qua trong phân tích độ phức tạp. Ta có thể ước
lượng tiếp bởi
\[
  O\left(
  \sum_{\substack{p\mid m,\ p\in\mathbb P}}
  \sum_{\substack{q\in\mathbb P,\ e\ge 1\\q^e\Vert m,\ p\ne q}}
  e\cdot d\left(\frac{m}{pq^e}\right)
  \right).
\]
Do đó có thể phân tích bằng tổng
\[
\begin{aligned}
&\sum_{m\in[n-L,n]}
  \sum_{\substack{p<n^{1/l}\\p\in\mathbb P,\ p\mid m}}
  \sum_{\substack{q<n^{1/l}\\q\in\mathbb P,\ e\ge 1,\ p\ne q}}
  d\left(\frac{m}{pq^e}\right)\\
&\le
  \sum_{\substack{p\le n^{1/l}\\p\in\mathbb P}}
  \sum_{\substack{q\le n^{1/l}\\q\in\mathbb P,\ e\ge 1}}
  \sum_{i\in\left[\frac{n-L}{pq^e},\,\frac{n}{pq^e}\right]} d(i).
\end{aligned}
\]

Chú ý rằng \(p,q,e\) chỉ có \(O(n^{2/l}/\ln n)\) khả năng, nên tiếp tục dùng
ước lượng tổng tiền tố hàm ước trong số học giải tích \([1]\):
\[
  \sum_{i\in[L,R]}d(i)
  =
  (R-L+1)\log R+O(R^{131/416}).
\]
Chỉ cần hằng số \(l\) thỏa
\[
  \frac{2}{l}+\frac{131}{416}<\frac12,
\]
ta có thể bỏ qua \(O(R^{131/416})\).

Vì vậy có thể ước lượng độ phức tạp tiếp thành
\[
  \sum_{\substack{p\le n^{1/l}\\p\in\mathbb P}}
  \sum_{\substack{q\le n^{1/l}\\q\in\mathbb P,\ e\ge 1}}
  \frac{L}{pq^e}.
\]
Do
\[
  \sum_{\substack{p\le S\\p\in\mathbb P}}\frac1p
  = O(\log\log S),
\]
và
\[
  \sum_{\substack{p\le S\\p\in\mathbb P,\ e\ge 1}}\frac1{p^e}
\]
không vượt quá hai lần tổng trước, phần độ phức tạp này ước lượng được là
\[
  L(\log\log n)^2.
\]

Tổng độ phức tạp là
\[
  \frac{n\log^2 n}{S}+S(\log\log n)^2.
\]
Khi lấy
\[
  S=\frac{\sqrt n\log n}{\log\log n},
\]
độ phức tạp thời gian là
\[
  O(\sqrt n\log n\log\log n).
\]

\subsection*{5. Tổng kết}

Trong bài viết này, tác giả thu được một phương pháp sàng có tính mở rộng khá
tốt, đưa tích chập sàng khối, phép chia sàng khối, tổng tiền tố sàng khối của
hàm nhân tính và bài toán thống kê tiền tố số nguyên tố về độ phức tạp thời
gian \(O(\sqrt n\operatorname{polylog} n)\). Hai bài toán sau cần thỏa một số
điều kiện, chẳng hạn \(f(p)\) là đa thức bậc hằng theo \(p\), hoặc một vài tính
chất đặc biệt khác. Bài viết cũng thảo luận tối ưu phương pháp này, đưa tích
chập sàng khối và phép chia sàng khối của hàm số học về
\[
  O(\sqrt n\log n\log\log n),
\]
đồng thời đưa tổng tiền tố sàng khối của hàm nhân tính và thống kê tiền tố số
nguyên tố về
\[
  O\left(\frac{\sqrt n\log^2 n}{\sqrt{\log\log n}}\right).
\]

Tuy độ phức tạp thời gian tốt hơn, do hằng số lớn và có nhiều thừa số \(\log\),
cách làm này trong phạm vi dữ liệu OI hiện nay vẫn không thể vượt qua sàng
Min\_25. Tác giả hy vọng người đọc tiếp tục suy nghĩ và nghiên cứu bài toán
tính tổng hàm nhân tính, và mong sẽ xuất hiện các cách làm thực dụng hơn.

\subsection*{6. Lời cảm ơn}

Cảm ơn China Computer Federation đã cung cấp nền tảng học tập và giao lưu.

Cảm ơn các huấn luyện viên đội tuyển tập huấn quốc gia, Peng Sijin và Yang
Yaoliang, đã hướng dẫn.

Cảm ơn các thầy ở Trường Trung học Học Quân, trong đó có Xu Xianyou, đã quan
tâm và hướng dẫn.

Cảm ơn gia đình đã đồng hành và ủng hộ.

Cảm ơn các bạn học và anh chị trong đội bạn đã giúp đỡ và đồng hành.

Cảm ơn các bạn học Zhang Yongxuan, Zhang Hengyi và những người khác đã đọc bản
thảo và góp ý cho bài viết này.

Cảm ơn các thầy cô và bạn học khác đã giúp đỡ tác giả.

\subsection*{Tài liệu tham khảo}

\begin{enumerate}
  \item \href{https://en.wikipedia.org/wiki/Divisor_summatory_function}{Divisor
  summatory function}.
  \item Li Baitian. \emph{Khung lý thuyết tính toán hàm sinh trong thi tin
  học}. Tập luận văn đội tuyển quốc gia Trung Quốc Olympic Tin học năm 2021,
  2021.
\end{enumerate}

\section{Chủ đề chính}

Từ mục lục, tập năm 2024 bao phủ các nhóm chủ đề sau:
\begin{itemize}
  \item hàm nhân tính, số học, liên phân số, thuật toán Euclid và \(q\)-analogue;
  \item matroid tuyến tính, matching có trọng số đỉnh và tô màu danh sách;
  \item cấu trúc dữ liệu, cập nhật-truy vấn đoạn và các bài toán trên cây;
  \item quy hoạch động, hợp nhất thông tin, bao hàm--loại trừ và nghịch đảo;
  \item hàm sinh, quá trình ngẫu nhiên, đếm đường đi trên lưới;
  \item tối ưu hằng số trên phần cứng hiện đại, lý thuyết tính toán và độ phức tạp;
  \item lattice, phân tích đa thức hệ số nguyên, chu trình tỉ lệ nhỏ nhất và các báo cáo ra đề.
\end{itemize}

\section{Ghi chú về OCR}

Phần mục lục của tập 2024 trích xuất được rõ ràng, nhưng lớp văn bản thân bài
không ổn định: nhiều đoạn trở thành mojibake và công cụ PDF báo cảnh báo phông
chữ. Bài đầu tiên ở trên được dịch từ OCR tại \texttt{/tmp/oipl-ocr/2024-p9-18.txt}
và đối chiếu với ảnh render các trang PDF 9--18. Một số công thức dài trong
phần 4 được giữ ở mức có thể xác nhận từ ảnh trang; các biến phụ như \(\omega(m)\)
được hiểu theo ký hiệu chuẩn là số lượng thừa số nguyên tố phân biệt.

\end{document}
